\documentclass{article}
\usepackage[final]{neurips_2026}
\usepackage[utf8]{inputenc}
\usepackage[T1]{fontenc}
\usepackage{hyperref}
\usepackage{url}
\usepackage{booktabs}
\usepackage{amsfonts}
\usepackage{amsmath}
\usepackage{nicefrac}
\usepackage{microtype}
\usepackage{xcolor}
\usepackage{graphicx}
\usepackage{enumitem}

\title{Lessons from Sim-to-Real on a 15-DOF Bipedal Robot:\\Hardware Realities, Reward Brittleness, and the Limits of BC Diffusion for Locomotion}

\author{%
  Zhenghao Xu\\
  Elmore Family School of Electrical and Computer Engineering\\
  Purdue University\\
  West Lafayette, IN 47906 \\
  \texttt{xu1783@purdue.edu} \\
}

\begin{document}

\maketitle

\begin{abstract}
We present a reflective account of working with the Open Duck Mini V2, a low-cost 15-DOF bipedal robot. The work spanned hardware assembly, PPO policy training in MuJoCo MJX, sim-to-real transfer attempts, behavior cloning of a Diffusion Policy from PPO rollouts, and on-policy refinement via DPPO. The central outcomes are largely negative: a PPO policy that dominates in simulation (beating the author-provided baseline on 11/14 metrics) fails to transfer to hardware; a BC diffusion policy fails catastrophically even in simulation (100\% fall rate); and DPPO, while tripling survival time over BC, still fails to produce a stable walker (100\% fall rate). Rather than presenting selective successes, we analyze where each stage broke down. We identify root causes in hardware assembly errors, reward overfitting in simulation, and the difficulty of adapting diffusion-based methods to bipedal locomotion. We conclude with concrete recommendations---system identification, real-world data collection, and curriculum training---for making progress on this class of problems.
\end{abstract}


\section{Introduction}

Sim-to-real transfer remains one of the central challenges in robot learning. The promise is compelling: train policies endlessly in simulation at zero marginal cost, then deploy on hardware. In practice, the gap between simulated and real dynamics---the \emph{sim-to-real gap}---makes this transfer fragile, especially for inherently unstable systems like bipedal robots.

This paper recounts a hands-on project with the Open Duck Mini V2, a small, low-cost ($\sim$\$200) bipedal robot with 15 degrees of freedom. The work progressed through five stages:

\begin{enumerate}[nosep]
  \item \textbf{Hardware assembly} of the physical robot from a kit
  \item \textbf{PPO policy training} in MuJoCo MJX, iterated over 28 experiments
  \item \textbf{Sim-to-real transfer} attempts with the trained policies
  \item \textbf{Diffusion Policy} via behavior cloning from PPO rollouts
  \item \textbf{DPPO} fine-tuning of the BC diffusion policy with on-policy RL
\end{enumerate}

Each stage produced a negative result or an instructive failure. The trained PPO policy beat the original author's baseline on 11 of 14 simulation metrics but performed poorly on hardware. The BC diffusion policy achieved a 100\% fall rate in simulation. DPPO improved over BC---surviving $2.5\times$ longer and accumulating $3\times$ more reward---but still fell in every episode. Rather than presenting selective successes, we document these failures honestly and analyze their root causes. We believe this kind of failure analysis is underrepresented in the literature and is often more instructive than success stories.

Our contributions are:
\begin{itemize}[nosep]
  \item A detailed analysis of how hardware assembly errors (missing Loctite, mis-initialized joints) create sim-to-real gaps that domain randomization cannot close
  \item A case study in reward overfitting: a policy that dominates in simulation but fails on hardware
  \item An empirical demonstration that BC-from-distillation fails catastrophically on bipedal locomotion, and that DPPO---while directionally correct---is insufficient in this regime without further engineering
  \item Concrete recommendations for future work: system identification, real2sim2real loops, and curriculum training
\end{itemize}


\section{Related Work}

\paragraph{Sim-to-real transfer.} Domain randomization \citep{tobin2017domain, peng2018sim} is the standard approach for bridging the sim-to-real gap, treating modeling errors as noise to be averaged over. System identification \citep{holladay2024robot} takes the complementary approach of directly measuring and matching real-world dynamics. Our experience suggests that for bipedal robots with significant hardware variability, SysID is essential and randomization alone is insufficient.

\paragraph{Diffusion Policy.} Diffusion Policy \citep{chi2023diffusion} applies denoising diffusion models to visuomotor policy learning, demonstrating strong performance on manipulation tasks. DPPO \citep{ren2024dppo} extends this by using a diffusion policy as the policy class within a PPO objective, combining diffusion's expressive action representation with RL's corrective feedback loop.

\paragraph{Behavior cloning limitations.} The compounding error problem in BC was formalized by \citet{ross2011dagger}. We discuss this in detail in Section~\ref{sec:bc_analysis} in the context of our empirical results.


\section{The Platform}

The Open Duck Mini V2 is a 15-DOF bipedal robot: 10 leg joints (hip yaw/roll/pitch, knee, ankle per leg) and 5 head/neck joints. It runs on a Raspberry Pi 5 with Feetech STS3215 servos and an onboard IMU. The upstream repository provides a complete PPO training pipeline using Brax in MuJoCo MJX, with automatic ONNX export and a runtime for the Pi.

The robot's appeal is its low cost and open-source toolchain. The challenge is that it is a biped---inherently less stable than a quadruped, and far more sensitive to sim-to-real gaps in joint dynamics, mass distribution, and sensor calibration.

The observation space is 101-dimensional: gyroscope (3), accelerometer (3), velocity command (7), joint angles minus defaults (14), joint velocities (14), three steps of action history ($3 \times 14$), motor targets (14), foot contacts (2), and gait phase (2). The action space is 14-dimensional: position offsets applied as $\text{motor\_targets} = \text{default} + \text{action} \times 0.25$. The policy runs at 50\,Hz (20\,ms control period) with physics at 500\,Hz.


\section{Hardware Lessons}

\subsection{Missing Loctite on Actuator Screws}

The assembly guide did not emphasize that threadlocker (Loctite) is essential for all screws on the actuator joints. After several sessions of testing, screws on multiple actuators loosened progressively, introducing mechanical play---backlash and friction that the simulated model did not account for.

The Feetech servos produce enough torque to vibrate screws loose over hundreds of gait cycles. The resulting backlash is asymmetric, velocity-dependent, and differs from joint to joint. The backlash simulation in MuJoCo (\texttt{scene\_flat\_terrain\_backlash.xml}) models generic gear play but cannot capture this specific pattern of screw-loosening-induced play.

\paragraph{Lesson.} Always use threadlocker on any joint that transmits force. This is a hardware problem that demands a hardware fix, not a software one.

\subsection{Incorrect Neck Joint Initialization}

The assembly guide did not provide clear instructions for initializing the neck joints (neck\_pitch, head\_pitch, head\_yaw, head\_roll) to their zero positions. We suspect the joints were initialized at an offset from where the simulated model expects them.

For a biped, this matters more than it might seem. The head assembly has non-trivial mass, and its position affects the center of mass of the entire robot. On a quadruped, a small COM offset is a nuisance. On a biped, it can be the difference between balancing and falling.

\paragraph{Lesson.} Before any sim-to-real deployment, verify every joint's zero position against the simulated model's home keyframe. The runtime's \texttt{init\_pos} array must exactly match \texttt{keyframe("home").ctrl} from the MJCF model.


\section{PPO Policy Training}
\label{sec:ppo}

\subsection{Training Pipeline}

We trained PPO policies using the MuJoCo MJX pipeline on rented cloud GPUs (Vast.ai, RTX 4060 Ti). Over 28 experiments, we iterated on reward weights through stacked fine-tuning: starting from a baseline policy, progressively adjusting reward scales and continuing training from the best checkpoint.

The key parameter discovery was that \texttt{tracking\_ang\_vel} dominates forward walking quality. Increasing it from the default 6.0 to 14.0 (through a chain of $6 \rightarrow 8 \rightarrow 10 \rightarrow 12 \rightarrow 14$) steadily improved angular velocity tracking and forward gait, at the cost of lateral tracking. The best policy (experiment 28, 159M steps) achieved the results shown in Table~\ref{tab:ppo_comparison}.

\begin{table}[t]
\caption{PPO policy comparison: custom policy (28 experiments of reward tuning) vs.\ the original author's best policy. The custom policy beats the baseline on 11/14 metrics in simulation.}
\label{tab:ppo_comparison}
\centering
\begin{tabular}{lccc}
\toprule
Metric & Custom Policy & Author Baseline & Delta \\
\midrule
Stand tracking (angular) & 0.907 & 0.883 & +2.7\% \\
Forward tracking (angular) & 0.655 & 0.520 & +26\% \\
Backward tracking (linear) & 0.229 & 0.073 & +214\% \\
Turn right (angular) & 0.526 & 0.286 & +84\% \\
Turn left (angular) & 0.547 & 0.322 & +70\% \\
\midrule
Metrics improved & \multicolumn{3}{c}{\textbf{11 / 14}} \\
\bottomrule
\end{tabular}
\end{table}


\subsection{The Sim-to-Real Gap}

Despite dominating in simulation, the custom policy performed poorly on the real robot. The robot could stand and hold position, but walking was unstable---it would fall after a few steps, especially during turns and lateral movement.

In retrospect, this is a classic case of \emph{reward overfitting in simulation}. By aggressively optimizing \texttt{tracking\_ang\_vel} through 28 experiments, we produced a policy that exploits the simulated dynamics to achieve high reward scores. The policy learned behaviors that are brittle: they work in the narrow band of dynamics the simulator provides but fail when the real dynamics diverge.

The original author's policy, while scoring lower in simulation, was validated on real hardware during development. It learned a more conservative gait that is robust to the sim-to-real gap. Our policy learned a more aggressive gait that scores higher in sim but lacks this robustness.

\paragraph{Lesson.} Beating an author-provided baseline in simulation does not mean you have a better policy for the real robot. If your goal is sim-to-real transfer, evaluate on hardware early and often. Simulation metrics are necessary but not sufficient.


\subsection{Domain Randomization Was Not Enough}

Our primary strategy for sim-to-real transfer was domain randomization during training: varying floor friction, joint PID gains, mass distribution, initial poses, and applying random pushes.

The problem is that domain randomization treats the sim-to-real gap as \emph{zero-mean noise} around the simulated parameters. It assumes the real robot's dynamics are somewhere within the randomization envelope. But when the real robot has systematic biases---loosened screws creating asymmetric backlash, a shifted center of mass, servo responses that differ from the identified model---no amount of randomization centered on the wrong parameters will help. You cannot randomize your way out of a systematic modeling error.


\section{Diffusion Policy via Behavior Cloning}

\subsection{Motivation and Setup}

We applied Diffusion Policy \citep{chi2023diffusion} to this platform via behavior cloning from PPO rollouts: generate $(o_t, a_t)$ pairs by rolling out the PPO policy in simulation, then train a Conditional U-Net to predict action chunks via denoising.

\begin{itemize}[nosep]
  \item \textbf{Dataset:} 4,755 trajectories from the original author's PPO policy, collected at 50\,Hz with randomized velocity commands, friction, initial poses, and impulse perturbations. Total: 865,780 steps.
  \item \textbf{Model:} ConditionalUnet1D (16.4M parameters) with DDIM sampling (16 denoising steps). Action chunk size 16, observation horizon 2.
  \item \textbf{Training:} 100 epochs, batch size 256, AdamW with cosine LR. Final validation loss: 0.0221.
\end{itemize}


\subsection{Results}

The diffusion policy failed catastrophically. Table~\ref{tab:diffusion_results} shows the head-to-head comparison.

\begin{table}[t]
\caption{Sim evaluation: PPO baseline vs.\ BC diffusion policy. The diffusion policy achieves 98\% fall rate and survives only 9 steps on average.}
\label{tab:diffusion_results}
\centering
\begin{tabular}{lccc}
\toprule
Metric & PPO Baseline & Diffusion BC & Delta \\
\midrule
Avg survival steps & 100 / 1000 & 9 / 1000 & $-$91\% \\
Fall rate & 0\% & 98\% & +98\,pp \\
Stand tracking (linear) & 0.809 & 0.125 & $-$85\% \\
Action smoothness (rate) & 0.053 & 0.357 & $6.7\times$ noisier \\
\bottomrule
\end{tabular}
\end{table}

Only 20\% of episodes survived even the \texttt{stand} command. All motion commands caused immediate falls.


\subsection{Analysis: Why BC Diffusion Fails on Bipedal Locomotion}
\label{sec:bc_analysis}

\paragraph{Compounding error (distribution shift).}
BC trains a mapping from observations to actions using expert data. At test time, any small prediction error causes the state to leave the training distribution. The next observation is now out-of-distribution, producing a worse action, compounding until the robot falls. This is the DAgger problem \citep{ross2011dagger}: BC policies are correct only on states visited by the expert, but must act on states produced by their own imperfect predictions.

\paragraph{Action smoothness.}
The diffusion policy's action rate was 5--50$\times$ higher than PPO's. Even with DDIM's deterministic sampling, the denoised action chunks had high inter-step variance. For a bipedal robot that requires smooth, coordinated joint trajectories, jerky actions are immediately destabilizing.

\paragraph{Unimodal teacher.}
Distilling from a single PPO policy produces a unimodal action distribution. The diffusion model only sees one way to respond to each situation and cannot learn the diverse recovery behaviors that PPO implicitly uses to stay upright. PPO developed these through the RL survival reward; the diffusion policy never experienced this pressure.

\paragraph{Locomotion vs.\ manipulation.}
This failure highlights a fundamental distinction between locomotion and manipulation. Diffusion Policy was demonstrated on manipulation tasks where small action errors are recoverable---a slightly misaligned gripper can be corrected on the next step. In bipedal locomotion, the state space is far less forgiving: a small error in foot placement can cause an unrecoverable fall within one or two steps. This is precisely why DPPO \citep{ren2024dppo}, which combines diffusion's expressive action representation with RL's corrective feedback, is the natural next step for locomotion.


\section{DPPO: On-Policy Refinement of the BC Diffusion Policy}
\label{sec:dppo}

\subsection{Motivation}

The BC diffusion policy's failure is attributable to distribution shift and the lack of corrective feedback. DPPO \citep{ren2024dppo} addresses this by using the diffusion model as the policy class within an on-policy PPO objective: the policy collects rollouts in the environment, receives reward signals, and updates via policy gradient. This provides the corrective feedback that BC lacks.

\subsection{Setup}

We initialized the DPPO policy from the BC diffusion checkpoint and trained for 500 iterations using the DPPO reference implementation adapted to the Open Duck environment.

\begin{itemize}[nosep]
  \item \textbf{Initialization:} BC diffusion policy (ConditionalUnet1D, converted to DPPO format)
  \item \textbf{Training:} 500 iterations, 20 parallel environments, 200 steps per iteration ($\sim$8 min/iteration on Vast.ai RTX 3090)
  \item \textbf{Domain randomization:} random pushes (0.3--1.5\,N every 80--160 steps), random friction ($0.5\text{--}1.5\times$), random initial joint noise, random yaw
  \item \textbf{Wrappers:} observation/action normalization, multi-step wrapper (2 obs steps, 4 action steps)
  \item \textbf{Evaluation:} 10 episodes per checkpoint, max 1000 steps (20\,s at 50\,Hz), seed\,=\,42
\end{itemize}

\subsection{Results}

Both BC and DPPO fall in every episode (100\% fall rate). However, DPPO produces meaningful improvement in survival time and cumulative reward, as shown in Table~\ref{tab:dppo_results}.

\begin{table}[t]
\caption{Sim evaluation: BC baseline vs.\ DPPO (best checkpoint, iteration 300). Both policies fall in every episode, but DPPO survives $2.5\times$ longer and accumulates $3\times$ more reward.}
\label{tab:dppo_results}
\centering
\begin{tabular}{lccc}
\toprule
Metric & BC Baseline & DPPO (itr 300) & Delta \\
\midrule
Fall rate & 100\% & 100\% & $-$ \\
Avg reward & $1368 \pm 224$ & $3992 \pm 2189$ & $+2.9\times$ \\
Avg survival (env steps) & $\sim$56 & $\sim$144 & $+2.6\times$ \\
Best episode reward & 1918 & 8191 & $+4.3\times$ \\
\bottomrule
\end{tabular}
\end{table}

The best checkpoint (iteration 300) achieved eval reward of 4009, up from the BC baseline of $\sim$3009. Later iterations showed degradation (3100--3810 range), suggesting overtraining.

\subsection{Analysis}

\paragraph{DPPO helps but is not enough.}
The $3\times$ reward improvement and $2.6\times$ survival improvement confirm that on-policy RL feedback addresses some of BC's weaknesses. The policy learns to delay falls by producing more stable initial actions. However, $2.9$\,s of average survival in a 20\,s episode is far from a working walker. The core difficulty remains: bipedal balance is a narrow-waist problem where compounding errors quickly lead to unrecoverable states, and 500 iterations of DPPO do not provide enough exploration to learn robust recovery behaviors.

\paragraph{High variance.}
The reward standard deviation for DPPO ($\pm 2189$, or $\sim$55\% of the mean) is much larger than for BC ($\pm 224$, or $\sim$16\%). This suggests the policy is highly sensitive to initial conditions and perturbations---it sometimes finds a stable trajectory briefly, but cannot consistently reproduce it.

\paragraph{Training instability.}
The degradation after iteration 300 (from 4009 to 3100--3810) is consistent with the reward overfitting pattern observed in PPO training (Section~\ref{sec:ppo}): the policy overfits to the training distribution and loses robustness. For DPPO, this may be exacerbated by the diffusion policy's high-dimensional action space making the optimization landscape more rugged.

\paragraph{Action chunking mismatch.}
The DPPO implementation uses 4 action steps per multi-step, while the original BC policy was trained with 16-step chunks. This architectural mismatch may limit how well the BC initialization transfers to the DPPO fine-tuning regime. A matched chunk size could improve results.


\section{What We Would Do Differently}

\subsection{System Identification Before Training}

The single most impactful change would be system identification (SysID) on the real robot before training. Instead of broad domain randomization centered on generic parameters, we would:

\begin{enumerate}[nosep]
  \item Measure each joint's actual response to commanded positions at various speeds and loads
  \item Identify the actual backlash profile per joint (likely asymmetric and joint-specific after screw loosening)
  \item Measure the real robot's mass and COM by physical weighing
  \item Update the MuJoCo model's joint parameters, friction coefficients, and body masses to match
\end{enumerate}

Domain randomization is a workaround for an inaccurate model; SysID fixes the model.

\subsection{Real-World Data Collection and Fine-Tuning}

With more time, we would set up a data collection pipeline on the real robot:

\begin{enumerate}[nosep]
  \item Run the author's best policy on the real robot while recording observations and actions at 50\,Hz
  \item Use this real-world data to fine-tune via offline RL or to identify systematic observation--action mismatches
  \item Iterate: collect data $\rightarrow$ update model $\rightarrow$ retrain $\rightarrow$ evaluate $\rightarrow$ repeat
\end{enumerate}

This is the ``real2sim2real'' loop that has become standard in modern robot learning. Even a small amount of real-world data is more valuable for closing the sim-to-real gap than any amount of simulated data.

\subsection{Curriculum Training for Diffusion Policies}

Our DPPO results (Section~\ref{sec:dppo}) show that naive on-policy refinement from a BC initialization is insufficient for bipedal locomotion. A more promising approach would be curriculum training: start without domain randomization to let the policy learn basic balance, then gradually introduce perturbations and velocity commands. This mirrors the curriculum used in standard PPO training, where the agent first learns to stand before learning to walk.


\section{Conclusion}

This project produced three negative results: a PPO policy that excels in simulation but fails on hardware, a BC diffusion policy that fails even in simulation, and a DPPO policy that improves over BC but still cannot walk. All three failures are informative.

The sim-to-real gap was dominated by hardware issues (loose screws, mis-initialized joints) that domain randomization cannot fix. The correct response to a systematic sim-to-real gap is not more randomization but better modeling---system identification to bring the simulated dynamics closer to reality.

The diffusion policy experiments---BC and DPPO---paint a sobering picture of applying diffusion-based methods to bipedal locomotion. BC fails due to distribution shift and lack of corrective feedback. DPPO provides that feedback and improves survival by $2.5\times$, but remains far from producing a stable walker. The difficulty is not just the policy class or the training objective; it is the inherent instability of bipedal balance, where small errors compound rapidly into unrecoverable falls.

The most valuable outcome of this project is not any single result but the understanding of \emph{where} each method breaks down and \emph{why}. For future work, the path forward requires: SysID to close the sim-to-real gap, curriculum training to progressively build difficulty for diffusion policies, and real2sim2real loops to ground simulated training in real-world dynamics.


\section{Outlook}

This project concluded earlier than intended due to time constraints. However, the lessons learned here directly inform ongoing and planned work. The immediate next step is a collaborative project on hierarchical multi-agent policies for a team of four Open Duck robots, where the sim-to-real insights and policy-class comparisons from this warm-up project will be applied in a multi-agent setting. Beyond this, the core takeaways---that system identification matters more than domain randomization, that closed-loop feedback is essential for unstable systems, and that honest failure analysis accelerates learning---transfer across embodiments to upcoming work on quadrupedal locomotion and bimanual manipulation.


\begin{ack}
We thank the Open Duck Mini community and the original author (Antoine Pirrone) for the open-source hardware design, training pipeline, and reference policies.
\end{ack}


\bibliographystyle{plain}
\bibliography{references}

\end{document}
